\chapter{Timed Models}

\section{Core assumption}
\begin{defbox}
A timed model adds explicit timing information to distributed behavior, often through deadlines, timestamps, or bounded-delay constraints that are weaker than full synchrony but stronger than pure asynchrony.
\end{defbox}

\section{Typical properties}
\begin{itemize}
  \item Time is part of the specification.
  \item Correctness may depend on meeting deadlines.
  \item Timing constraints can be partial rather than global.
\end{itemize}

\section{Why it matters}
\begin{keybox}
Timed models are useful when real systems need temporal guarantees but cannot assume every operation is perfectly synchronized.
\end{keybox}

\section{Exam prompts}
\begin{itemize}
  \item Distinguish timed models from synchronous models.
  \item Explain why deadlines are central in this model.
  \item Describe one application where timing is a correctness concern.
\end{itemize}
